%! Author = sriadityavedantam
%! Date = 3/25/26

\lesson{48}{Self-Study Four}{Day 48}

\begin{definition}[Completion]
    Let $(M,d)$ be a metric space.
    Then there exists a metric space $(\overline{M},\overline{d})$ such that:
    \begin{enumerate}
        \item $M\subseteq\overline{M}$.
        \item $\overline{d}|_{M\times M} = d$.
        \item $\overline{M}$ is Cauchy complete.
        \item $\overline{M}$ is the closure of $M$.
    \end{enumerate}
\end{definition}

\begin{definition}[$C_\infty(M)$]
    \[
        C_\infty(M) \coloneqq \left\{f:M\to\rr : \text{$f$ is continuous and } \sup_{m\in M}\abs{f(m)}<\infty\right\}.
    \]
    This is a metric space with
    \[
        d_\infty(f,g) \coloneqq \sup_{m\in M}\abs{f(m)-g(m)}.
    \]
\end{definition}

\begin{definition}[Banach Space]
    A Banach space is a normed vector space that is Cauchy complete with respect to the metric induced by the norm.
    Examples include $\rr^n$, $\c^n$, $C_\infty(M)$, and $C^0([a,b])$ with the sup norm.
\end{definition}

\begin{definition}[Functional]
    A functional is a linear map $T:X\to\rr$ or $T:X\to\c$.
\end{definition}

\begin{definition}[Operator Norm]
    \[
        \norm{T}_{op} \coloneqq \sup_{\substack{x\in X\\ \norm{x} = 1}}\abs{T(x)}.
    \]
\end{definition}

\begin{remark}
    The space of bounded linear functionals is a Banach space.
\end{remark}

\begin{definition}[Inner Product Space]
    An inner product space is a vector space $X$ with a map
    \[
        \angled{\cdot,\cdot}:X\times X\to \rr
    \]
    such that:
    \begin{itemize}
        \item $\angled{x,y} = \angled{y,x}$,
        \item $\angled{ax + by,z} = a\angled{x,z} + b\angled{y,z}$,
        \item if $x\neq 0$, then $\angled{x,x} > 0$.
    \end{itemize}
\end{definition}

\begin{definition}[Hilbert Space]
    A Hilbert space is an inner product space that is complete with respect to the norm induced by the inner product.
\end{definition}

\begin{remark}
    An example where Riemann integration fails, and hence motivates measure theory, is the Dirichlet function:
    \[
        \mathbb{1}_{\q}(x) =
        \begin{cases}
            1, & x\in\q,\\
            0, & \text{otherwise.}
        \end{cases}
    \]
    This function is not Riemann integrable on any interval.
\end{remark}